\documentclass[10pt,a4paper]{article}
\usepackage[spanish,es-nodecimaldot]{babel}
\usepackage[margin=20mm]{geometry}
\usepackage{microtype,booktabs,tabularx,xcolor,fancyhdr}
\usepackage{pgfplots}
\usepackage[colorlinks=true,allcolors=azul]{hyperref}
\pgfplotsset{compat=1.18}

\definecolor{azul}{HTML}{163F6A}
\definecolor{verde}{HTML}{007960}
\definecolor{tinta}{HTML}{172B39}
\definecolor{veladura}{HTML}{EFF5F6}
\color{tinta}
\setlength{\parindent}{0pt}
\setlength{\parskip}{5pt}
\setlength{\headheight}{14pt}
\pagestyle{fancy}
\fancyhf{}
\fancyhead[L]{\small\textcolor{azul}{EDGE-DB · BENCHMARK EN RASPBERRY PI}}
\fancyhead[R]{\small Luis Alberto Nafate Hernandez}
\fancyfoot[L]{\small Caso de estudio · noviembre de 2026}
\fancyfoot[R]{\small\thepage}
\renewcommand{\headrulewidth}{0.3pt}
\newcommand{\entrada}[1]{\vspace{4pt}{\large\bfseries\color{azul}#1}\par}
\newcommand{\clave}[1]{\texttt{#1}}
% Generado por results/analizar.py a partir del ensayo indicado.
\newcommand{\EdgeDato}[1]{\csname edge-#1\endcsname}
\expandafter\def\csname edge-pilot_runs\endcsname{6}
\expandafter\def\csname edge-final_runs\endcsname{70}
\expandafter\def\csname edge-attempts\endcsname{80}
\expandafter\def\csname edge-failed\endcsname{1}
\expandafter\def\csname edge-superseded_completed\endcsname{3}
\expandafter\def\csname edge-low_p99\endcsname{64}
\expandafter\def\csname edge-total_operations\endcsname{2750}
\expandafter\def\csname edge-idle_samples\endcsname{62}
\expandafter\def\csname edge-idle_duration\endcsname{64.3}
\expandafter\def\csname edge-idle_ram\endcsname{256.1}
\expandafter\def\csname edge-idle_rss\endcsname{73.7}
\expandafter\def\csname edge-run_ram_min\endcsname{188.6}
\expandafter\def\csname edge-run_rss_max\endcsname{368.8}
\expandafter\def\csname edge-run_temp_min\endcsname{48.3}
\expandafter\def\csname edge-run_temp_max\endcsname{68.8}
\expandafter\def\csname edge-run_cpu_max\endcsname{99.8}
\expandafter\def\csname edge-run_swap_max\endcsname{70.2}
\expandafter\def\csname edge-run_read_max\endcsname{7303168}
\expandafter\def\csname edge-run_write_max\endcsname{93696000}
\expandafter\def\csname edge-trace_samples\endcsname{719}
\expandafter\def\csname edge-boot_check\endcsname{coincidió}
\expandafter\def\csname edge-kernel_oom\endcsname{sin coincidencias}
\expandafter\def\csname edge-flags\endcsname{0x0}
\expandafter\def\csname edge-window\endcsname{5}
\expandafter\def\csname edge-warm\endcsname{1}
\expandafter\def\csname edge-cool\endcsname{1}
\expandafter\def\csname edge-read_mib\endcsname{7.0}
\expandafter\def\csname edge-write_mib\endcsname{89.4}
\expandafter\def\csname edge-q1_1000_1_p95\endcsname{85.6}
\expandafter\def\csname edge-q1_1000_1_qps\endcsname{14.0}
\expandafter\def\csname edge-q1_10000_1_p95\endcsname{539.8}
\expandafter\def\csname edge-q1_10000_1_qps\endcsname{1.8}
\expandafter\def\csname edge-q1_10000_2_p95\endcsname{552.6}
\expandafter\def\csname edge-q1_10000_2_qps\endcsname{3.6}
\expandafter\def\csname edge-q1_10000_5_p95\endcsname{797.3}
\expandafter\def\csname edge-q1_10000_5_qps\endcsname{7.6}
\expandafter\def\csname edge-q1_50000_1_p95\endcsname{3115.0}
\expandafter\def\csname edge-q1_50000_1_qps\endcsname{0.3}
\expandafter\def\csname edge-q2_1000_1_p95\endcsname{160.6}
\expandafter\def\csname edge-q2_1000_1_qps\endcsname{20.4}
\expandafter\def\csname edge-q2_10000_1_p95\endcsname{220.6}
\expandafter\def\csname edge-q2_10000_1_qps\endcsname{6.4}
\expandafter\def\csname edge-q2_10000_2_p95\endcsname{246.9}
\expandafter\def\csname edge-q2_10000_2_qps\endcsname{11.6}
\expandafter\def\csname edge-q2_10000_5_p95\endcsname{473.5}
\expandafter\def\csname edge-q2_10000_5_qps\endcsname{21.2}
\expandafter\def\csname edge-q2_50000_1_p95\endcsname{313.3}
\expandafter\def\csname edge-q2_50000_1_qps\endcsname{3.4}
\expandafter\def\csname edge-50000-q1-nodes\endcsname{Aggregate, Incremental Sort, Nested Loop, Gather Merge, Sort, Seq Scan, Function Scan}
\expandafter\def\csname edge-50000-q1-sort\endcsname{external merge}
\expandafter\def\csname edge-50000-q1-temp\endcsname{3532}
\expandafter\def\csname edge-50000-q1-sharedread\endcsname{0}
\expandafter\def\csname edge-50000-q1-sharedhit\endcsname{3938}
\expandafter\def\csname edge-50000-q1-ms\endcsname{3179.7}
\expandafter\def\csname edge-50000-q2-nodes\endcsname{Limit, Gather Merge, Sort, Seq Scan}
\expandafter\def\csname edge-50000-q2-sort\endcsname{top-N heapsort}
\expandafter\def\csname edge-50000-q2-temp\endcsname{0}
\expandafter\def\csname edge-50000-q2-sharedread\endcsname{0}
\expandafter\def\csname edge-50000-q2-sharedhit\endcsname{3950}
\expandafter\def\csname edge-50000-q2-ms\endcsname{305.6}
\expandafter\def\csname edge-50000-q1-jit-off-ms\endcsname{3057.4}
\expandafter\def\csname edge-50000-q1-jit-functions\endcsname{17}
\newcommand{\EdgeRows}{%
1,000 & 1 & Q1 & 62 & 85.6 [82.3, 266.3] & 290.3 & 14.0 \\
1,000 & 1 & Q2 & 84 & 160.6 [57.8, 216.3] & 304.7 & 20.4 \\
5,000 & 1 & Q1 & 14 & 469.3 [252.8, 501.5] & 469.3 & 3.6 \\
5,000 & 1 & Q2 & 10 & 115.3 [112.8, 134.2] & 321.3 & 10.8 \\
10,000 & 1 & Q1 & 9 & 539.8 [505.6, 540.4] & 539.8 & 1.8 \\
10,000 & 1 & Q2 & 25 & 220.6 [199.2, 395.3] & 397.3 & 6.4 \\
10,000 & 2 & Q1 & 18 & 552.6 [541.5, 565.9] & 552.6 & 3.6 \\
10,000 & 2 & Q2 & 44 & 246.9 [226.3, 269.5] & 302.1 & 11.6 \\
10,000 & 5 & Q1 & 37 & 797.3 [776.7, 839.4] & 862.5 & 7.6 \\
10,000 & 5 & Q2 & 92 & 473.5 [335.1, 525.1] & 681.8 & 21.2 \\
25,000 & 1 & Q1 & 2 & 1479.4 [1476.3, 1482.9] & 1479.4 & 0.6 \\
25,000 & 1 & Q2 & 14 & 329.4 [329.0, 410.4] & 329.4 & 3.0 \\
50,000 & 1 & Q1 & 6 & 3115.0 [3033.2, 3450.3] & 3115.0 & 0.3 \\
50,000 & 1 & Q2 & 15 & 313.3 [289.5, 461.2] & 313.3 & 3.4 \\
}
\newcommand{\EdgeDatasetRows}{%
1,000 & 3,023 & 598 & 0.35 & 0.95 \\
5,000 & 14,949 & 3,004 & 1.77 & 3.53 \\
10,000 & 29,939 & 5,960 & 3.54 & 6.64 \\
25,000 & 74,769 & 14,972 & 8.86 & 15.46 \\
50,000 & 149,760 & 30,171 & 17.73 & 30.99 \\
}
\newcommand{\EdgeLatency}{%
\addplot[azul,thick,mark=*] coordinates {(1000,85.642) (5000,469.275) (10000,539.841) (25000,1479.439) (50000,3115.039)};
\addlegendentry{Q1}
\addplot[verde,thick,mark=square*] coordinates {(1000,160.624) (5000,115.311) (10000,220.581) (25000,329.356) (50000,313.254)};
\addlegendentry{Q2}
}
\newcommand{\EdgeThroughput}{%
\addplot[azul,thick,mark=*] coordinates {(1,1.800) (2,3.600) (5,7.600)};
\addlegendentry{Q1}
\addplot[verde,thick,mark=square*] coordinates {(1,6.400) (2,11.600) (5,21.200)};
\addlegendentry{Q2}
}
\newcommand{\EdgeTemperature}{%
\addplot[verde,thick] coordinates {(0.000,47.774) (0.087,47.774) (0.174,48.312) (0.261,48.850) (0.348,47.774) (0.436,48.312) (0.522,48.312) (0.614,48.312) (0.697,48.312) (0.784,48.312) (0.871,49.388) (0.958,48.312) (1.045,49.388) (1.132,49.388) (1.219,49.388) (1.306,48.312) (1.393,49.388) (1.480,49.388) (1.568,49.388) (1.654,49.388) (1.742,48.850) (1.829,49.388) (1.916,49.388) (2.023,48.312) (2.091,49.388) (2.178,49.388) (2.265,49.388) (2.352,49.388) (2.439,50.464) (2.527,49.388) (2.616,49.388) (2.702,49.926) (2.787,50.464) (2.874,49.388) (2.961,50.464) (3.048,50.464) (3.135,50.464) (3.222,50.464) (3.311,49.926) (3.396,50.464) (3.483,51.002) (3.570,50.464) (3.657,50.464) (3.744,51.002) (3.831,49.926) (3.918,51.002) (4.007,51.540) (4.092,51.002) (4.179,51.002) (4.266,51.540) (4.353,51.540) (4.440,51.002) (4.527,51.540) (4.614,51.540) (4.701,51.540) (4.788,51.540) (4.875,52.078) (4.962,52.078) (5.049,52.616) (5.136,52.616) (5.231,51.540) (5.311,51.540) (5.397,52.078) (5.486,51.540) (5.571,53.154) (5.658,52.616) (5.745,52.616) (5.833,53.692) (5.920,52.616) (6.007,52.616) (6.094,53.154) (6.181,54.230) (6.268,54.768) (6.355,53.154) (6.443,56.382) (6.533,55.306)};
\addplot[verde,thick] coordinates {(8.345,49.388) (8.432,52.616) (8.524,49.388) (8.606,51.002) (8.698,51.540) (8.798,52.616) (8.885,52.616) (8.973,53.154) (9.060,52.616) (9.147,53.154) (9.234,53.692) (9.321,54.768) (9.408,53.692) (9.496,55.844) (9.583,54.768) (9.671,56.920) (9.758,55.306) (9.845,55.844) (9.933,55.306) (10.021,55.844) (10.108,54.768) (10.195,56.382) (10.282,54.768) (10.372,54.768) (10.461,55.306) (10.548,55.306) (10.634,56.920) (10.726,56.920) (10.809,55.844) (10.896,56.382) (10.983,56.920) (11.071,56.920) (11.158,55.306) (11.246,58.534) (11.333,57.996) (11.453,57.458) (11.507,57.996) (11.595,57.458) (11.682,57.996) (11.769,57.996) (11.856,57.996) (11.950,57.996) (12.030,57.458) (12.117,58.534) (12.205,59.072) (12.291,57.996) (12.379,57.996) (12.466,58.534) (12.556,59.072) (12.640,59.072) (12.727,57.458) (12.816,62.300) (12.906,62.300) (12.995,62.300) (13.085,62.838) (13.172,64.452) (13.261,64.452) (13.350,64.452) (13.442,66.604) (13.530,65.528) (13.622,64.452) (13.707,66.604) (13.798,68.218) (13.887,66.604) (13.979,65.528) (14.066,68.218) (14.155,65.528) (14.255,67.142)};
}
\newcommand{\EdgeRam}{%
\addplot[verde,thick] coordinates {(0.000,251.492) (0.087,260.367) (0.174,257.648) (0.261,253.176) (0.348,257.230) (0.436,254.008) (0.522,247.562) (0.614,252.906) (0.697,256.039) (0.784,254.246) (0.871,252.113) (0.958,253.762) (1.045,258.059) (1.132,254.082) (1.219,258.426) (1.306,252.945) (1.393,244.422) (1.480,257.047) (1.568,257.484) (1.654,259.504) (1.742,257.789) (1.829,256.539) (1.916,258.422) (2.023,252.820) (2.091,255.840) (2.178,254.961) (2.265,253.949) (2.352,245.246) (2.439,255.090) (2.527,253.074) (2.616,247.523) (2.702,252.020) (2.787,253.051) (2.874,256.418) (2.961,251.227) (3.048,254.777) (3.135,257.508) (3.222,252.016) (3.311,247.676) (3.396,246.293) (3.483,244.844) (3.570,246.996) (3.657,247.574) (3.744,245.223) (3.831,235.734) (3.918,248.691) (4.007,249.535) (4.092,252.273) (4.179,250.059) (4.266,248.352) (4.353,249.297) (4.440,248.680) (4.527,243.727) (4.614,245.109) (4.701,235.457) (4.788,221.207) (4.875,209.148) (4.962,203.176) (5.049,219.852) (5.136,228.254) (5.231,235.816) (5.311,231.973) (5.397,222.441) (5.486,234.910) (5.571,213.016) (5.658,214.727) (5.745,225.070) (5.833,203.785) (5.920,220.285) (6.007,230.203) (6.094,209.594) (6.181,209.141) (6.268,164.113) (6.355,208.641) (6.443,204.305) (6.533,170.766)};
\addplot[verde,thick] coordinates {(8.345,199.969) (8.432,197.875) (8.524,197.289) (8.606,200.371) (8.698,204.207) (8.798,197.535) (8.885,211.137) (8.973,198.434) (9.060,222.516) (9.147,204.391) (9.234,207.945) (9.321,202.336) (9.408,231.398) (9.496,235.266) (9.583,225.980) (9.671,236.863) (9.758,199.574) (9.845,200.234) (9.933,200.398) (10.021,197.602) (10.108,244.066) (10.195,201.633) (10.282,226.281) (10.372,193.961) (10.461,199.680) (10.548,228.895) (10.634,230.352) (10.726,218.699) (10.809,207.727) (10.896,194.980) (10.983,196.719) (11.071,200.141) (11.158,231.402) (11.246,229.047) (11.333,216.609) (11.453,225.520) (11.507,207.902) (11.595,207.906) (11.682,217.723) (11.769,215.590) (11.856,226.152) (11.950,213.352) (12.030,205.730) (12.117,219.758) (12.205,208.504) (12.291,232.508) (12.379,218.078) (12.466,210.473) (12.556,224.629) (12.640,218.766) (12.727,221.484) (12.816,206.910) (12.906,192.277) (12.995,210.289) (13.085,203.922) (13.172,213.293) (13.261,197.730) (13.350,220.852) (13.442,205.793) (13.530,182.695) (13.622,209.133) (13.707,207.250) (13.798,210.633) (13.887,211.004) (13.979,233.004) (14.066,209.430) (14.155,211.090) (14.255,206.473)};
}
\newcommand{\EdgeFrequency}{%
\addplot[verde,thick] coordinates {(0.000,1000.000) (0.087,1000.000) (0.174,1000.000) (0.261,1000.000) (0.348,700.000) (0.436,1000.000) (0.522,700.000) (0.614,1000.000) (0.697,1000.000) (0.784,1000.000) (0.871,1000.000) (0.958,1000.000) (1.045,1000.000) (1.132,1000.000) (1.219,1000.000) (1.306,1000.000) (1.393,1000.000) (1.480,1000.000) (1.568,1000.000) (1.654,1000.000) (1.742,1000.000) (1.829,1000.000) (1.916,1000.000) (2.023,1000.000) (2.091,1000.000) (2.178,1000.000) (2.265,1000.000) (2.352,1000.000) (2.439,1000.000) (2.527,1000.000) (2.616,1000.000) (2.702,1000.000) (2.787,1000.000) (2.874,1000.000) (2.961,1000.000) (3.048,1000.000) (3.135,1000.000) (3.222,1000.000) (3.311,1000.000) (3.396,1000.000) (3.483,1000.000) (3.570,1000.000) (3.657,1000.000) (3.744,1000.000) (3.831,1000.000) (3.918,1000.000) (4.007,1000.000) (4.092,1000.000) (4.179,1000.000) (4.266,1000.000) (4.353,1000.000) (4.440,1000.000) (4.527,1000.000) (4.614,1000.000) (4.701,1000.000) (4.788,1000.000) (4.875,1000.000) (4.962,1000.000) (5.049,1000.000) (5.136,1000.000) (5.231,1000.000) (5.311,1000.000) (5.397,1000.000) (5.486,1000.000) (5.571,1000.000) (5.658,1000.000) (5.745,1000.000) (5.833,1000.000) (5.920,1000.000) (6.007,1000.000) (6.094,1000.000) (6.181,1000.000) (6.268,1000.000) (6.355,1000.000) (6.443,1000.000) (6.533,1000.000)};
\addplot[verde,thick] coordinates {(8.345,1000.000) (8.432,1000.000) (8.524,1000.000) (8.606,1000.000) (8.698,1000.000) (8.798,1000.000) (8.885,1000.000) (8.973,1000.000) (9.060,1000.000) (9.147,1000.000) (9.234,1000.000) (9.321,1000.000) (9.408,1000.000) (9.496,1000.000) (9.583,1000.000) (9.671,1000.000) (9.758,1000.000) (9.845,1000.000) (9.933,1000.000) (10.021,1000.000) (10.108,1000.000) (10.195,1000.000) (10.282,1000.000) (10.372,1000.000) (10.461,1000.000) (10.548,1000.000) (10.634,1000.000) (10.726,1000.000) (10.809,1000.000) (10.896,1000.000) (10.983,1000.000) (11.071,1000.000) (11.158,1000.000) (11.246,1000.000) (11.333,1000.000) (11.453,1000.000) (11.507,1000.000) (11.595,1000.000) (11.682,1000.000) (11.769,1000.000) (11.856,1000.000) (11.950,1000.000) (12.030,1000.000) (12.117,1000.000) (12.205,1000.000) (12.291,1000.000) (12.379,1000.000) (12.466,1000.000) (12.556,1000.000) (12.640,1000.000) (12.727,1000.000) (12.816,1000.000) (12.906,1000.000) (12.995,1000.000) (13.085,1000.000) (13.172,1000.000) (13.261,1000.000) (13.350,1000.000) (13.442,1000.000) (13.530,1000.000) (13.622,1000.000) (13.707,1000.000) (13.798,1000.000) (13.887,1000.000) (13.979,1000.000) (14.066,1000.000) (14.155,1000.000) (14.255,1000.000)};
}


% Generado por results/analizar_extension.py.
\newcommand{\PgPerformance}{%
1,000 & Q1 & 1 & 5/6 & 14.53 & 54.4 & 110.9 [85.8, 261.3] & 281.3 & 0 \\
1,000 & Q2 & 1 & 5/5 & 23.25 & 25.0 & 162.3 [86.9, 177.9] & 321.5 & 0 \\
5,000 & Q1 & 1 & 5/5 & 3.95 & 244.3 & 281.1 [262.0, 394.9] & 471.4 & 0 \\
5,000 & Q2 & 1 & 5/5 & 9.82 & 81.0 & 162.8 [121.7, 166.2] & 407.2 & 0 \\
10,000 & Q1 & 1 & 5/5 & 2.03 & 484.3 & 504.0 [503.2, 510.1] & 573.2 & 1 \\
10,000 & Q2 & 1 & 5/5 & 7.33 & 133.9 & 156.6 [154.7, 163.0] & 188.7 & 0 \\
10,000 & Q1 & 2 & 5/5 & 4.02 & 487.5 & 521.6 [505.7, 540.0] & 706.1 & 0 \\
10,000 & Q2 & 2 & 5/6 & 13.57 & 139.8 & 181.8 [177.8, 187.1] & 236.1 & 0 \\
10,000 & Q1 & 5 & 5/5 & 7.80 & 623.0 & 811.1 [796.5, 833.5] & 947.9 & 0 \\
10,000 & Q2 & 5 & 5/5 & 27.90 & 171.9 & 236.7 [230.0, 241.8] & 308.9 & 0 \\
25,000 & Q1 & 1 & 5/5 & 0.68 & 1453.6 & 1491.2 [1489.4, 1545.4] & 1556.3 & 5 \\
25,000 & Q2 & 1 & 5/5 & 3.37 & 286.2 & 352.3 [337.8, 366.8] & 400.0 & 0 \\
50,000 & Q1 & 1 & 5/5 & 0.33 & 2803.2 & 3133.3 [3073.8, 3136.9] & 3189.3 & 5 \\
50,000 & Q2 & 1 & 5/5 & 3.48 & 250.2 & 469.6 [456.5, 481.0] & 1075.4 & 0 \\
}
\newcommand{\PgResources}{%
1,000 & Q1 & 1 & 19.7 & 91.9 & 254.3 & 49.9 & 0.0 & 0.0 \\
1,000 & Q2 & 1 & 11.6 & 91.4 & 252.5 & 48.3 & 0.0 & 0.0 \\
5,000 & Q1 & 1 & 25.1 & 95.3 & 248.1 & 51.0 & 0.0 & 0.0 \\
5,000 & Q2 & 1 & 18.8 & 93.7 & 250.0 & 50.5 & 0.0 & 0.0 \\
10,000 & Q1 & 1 & 25.8 & 100.0 & 235.5 & 50.5 & 0.0 & 0.0 \\
10,000 & Q2 & 1 & 24.2 & 96.7 & 241.5 & 50.5 & 0.0 & 0.0 \\
10,000 & Q1 & 2 & 50.0 & 126.2 & 233.2 & 56.4 & 0.0 & 0.0 \\
10,000 & Q2 & 2 & 44.3 & 119.7 & 238.5 & 55.8 & 0.0 & 0.0 \\
10,000 & Q1 & 5 & 98.3 & 204.3 & 213.1 & 66.6 & 0.0 & 0.0 \\
10,000 & Q2 & 5 & 92.9 & 188.2 & 226.3 & 65.5 & 0.0 & 0.0 \\
25,000 & Q1 & 1 & 31.3 & 256.6 & 198.2 & 52.6 & 0.0 & 0.5 \\
25,000 & Q2 & 1 & 24.5 & 105.7 & 228.3 & 52.1 & 0.0 & 0.0 \\
50,000 & Q1 & 1 & 32.9 & 333.5 & 190.5 & 55.3 & 1.5 & 227.8 \\
50,000 & Q2 & 1 & 58.9 & 172.6 & 221.2 & 59.6 & 0.0 & 0.0 \\
}
\newcommand{\PgOneVolume}{%
(1000,110.9402) (5000,281.0834) (10000,503.9962) (25000,1491.1514) (50000,3133.2569)
}
\newcommand{\PgOneClients}{%
(1,2.0333) (2,4.0167) (5,7.8000)
}
\newcommand{\PgTwoVolume}{%
(1000,162.3468) (5000,162.8116) (10000,156.5683) (25000,352.3388) (50000,469.6297)
}
\newcommand{\PgTwoClients}{%
(1,7.3333) (2,13.5667) (5,27.9000)
}
\newcommand{\PgWrites}{%
227.8
}
\newcommand{\PgTradeoff}{%
(7.3333,156.568) [1] (13.5667,181.793) [2] (27.9000,236.657) [5]
}
\newcommand{\PgCooling}{%
(base,41) (extendida,24) (limite,3)
}
\newcommand{\PgCoolingEvaluable}{%
68
}
\newcommand{\PgCoolingProlonged}{%
27
}
\newcommand{\PgCoolingUnrecovered}{%
3
}
\newcommand{\MongoPerformance}{%
1,000 & Q1 & 1 & 5/5 & 11.98 & 76.5 & 97.8 [91.0, 103.4] & 315.4 & 0 \\
1,000 & Q2 & 1 & 5/6 & 10.37 & 76.8 & 190.8 [136.3, 225.3] & 393.5 & 0 \\
5,000 & Q1 & 1 & 5/5 & 2.80 & 345.3 & 377.7 [374.0, 409.8] & 630.5 & 0 \\
5,000 & Q2 & 1 & 5/5 & 3.92 & 206.7 & 483.0 [364.6, 592.7] & 1175.7 & 0 \\
10,000 & Q1 & 1 & 5/5 & 1.38 & 688.8 & 904.7 [760.6, 1130.3] & 1179.1 & 5 \\
10,000 & Q2 & 1 & 5/6 & 2.35 & 317.7 & 975.9 [650.3, 1214.8] & 2140.4 & 0 \\
10,000 & Q1 & 2 & 5/6 & 2.65 & 688.3 & 1030.5 [1014.9, 1046.3] & 1746.9 & 0 \\
10,000 & Q2 & 2 & 5/5 & 4.30 & 322.5 & 1384.1 [1068.8, 1537.8] & 2346.5 & 0 \\
10,000 & Q1 & 5 & 5/6 & 5.13 & 886.1 & 1558.8 [1454.8, 1712.1] & 2461.3 & 0 \\
10,000 & Q2 & 5 & 5/6 & 8.27 & 398.1 & 1719.7 [1062.9, 2271.7] & 3012.8 & 0 \\
25,000 & Q1 & 1 & 5/5 & 0.58 & 1675.8 & 1917.9 [1912.6, 1954.0] & 2089.7 & 5 \\
25,000 & Q2 & 1 & 5/7 & 1.62 & 523.9 & 1242.5 [1162.7, 1320.6] & 1803.7 & 4 \\
50,000 & Q1 & 1 & 5/6 & 0.28 & 3356.7 & 4148.1 [3810.0, 4266.2] & 4148.1 & 5 \\
50,000 & Q2 & 1 & 5/5 & 0.83 & 980.8 & 2163.0 [1783.3, 4019.2] & 3978.1 & 5 \\
}
\newcommand{\MongoResources}{%
1,000 & Q1 & 1 & 22.7 & 153.5 & 200.8 & 50.5 & 0.0 & 0.6 \\
1,000 & Q2 & 1 & 10.8 & 153.5 & 202.2 & 47.8 & 0.0 & 0.6 \\
5,000 & Q1 & 1 & 25.2 & 147.8 & 199.2 & 52.1 & 0.0 & 0.6 \\
5,000 & Q2 & 1 & 13.1 & 147.9 & 199.5 & 49.9 & 0.0 & 0.6 \\
10,000 & Q1 & 1 & 25.2 & 167.8 & 188.8 & 52.6 & 0.0 & 0.6 \\
10,000 & Q2 & 1 & 13.9 & 166.1 & 188.5 & 50.5 & 0.0 & 0.5 \\
10,000 & Q1 & 2 & 47.2 & 171.8 & 181.0 & 58.0 & 0.0 & 0.7 \\
10,000 & Q2 & 2 & 24.9 & 171.8 & 180.3 & 54.8 & 0.0 & 0.6 \\
10,000 & Q1 & 5 & 90.4 & 195.4 & 164.2 & 67.7 & 0.0 & 0.6 \\
10,000 & Q2 & 5 & 48.7 & 169.2 & 191.1 & 60.7 & 0.0 & 0.6 \\
25,000 & Q1 & 1 & 26.0 & 188.7 & 167.9 & 53.2 & 0.0 & 0.6 \\
25,000 & Q2 & 1 & 20.2 & 188.4 & 169.2 & 52.6 & 0.0 & 0.6 \\
50,000 & Q1 & 1 & 26.0 & 182.0 & 182.9 & 53.2 & 0.0 & 0.6 \\
50,000 & Q2 & 1 & 19.9 & 184.1 & 181.0 & 52.6 & 0.0 & 0.6 \\
}
\newcommand{\MongoOneVolume}{%
(1000,97.7870) (5000,377.7371) (10000,904.7442) (25000,1917.9124) (50000,4148.0959)
}
\newcommand{\MongoOneClients}{%
(1,1.3833) (2,2.6500) (5,5.1333)
}
\newcommand{\MongoTwoVolume}{%
(1000,190.8448) (5000,483.0174) (10000,975.8790) (25000,1242.5086) (50000,2163.0279)
}
\newcommand{\MongoTwoClients}{%
(1,2.3500) (2,4.3000) (5,8.2667)
}
\newcommand{\MongoWrites}{%
0.6
}
\newcommand{\MongoTradeoff}{%
(2.3500,975.879) [1] (4.3000,1384.069) [2] (8.2667,1719.740) [5]
}
\newcommand{\MongoCooling}{%
(base,56) (extendida,14) (limite,0)
}
\newcommand{\MongoCoolingEvaluable}{%
70
}
\newcommand{\MongoCoolingProlonged}{%
14
}
\newcommand{\MongoCoolingUnrecovered}{%
0
}
\newcommand{\BudgetRows}{%
PostgreSQL & 469.6 & 266.4--652.2 & 4/5 \\
MongoDB comunitario & 2163.0 & 1165.8--4770.0 & 0/5 \\
}
\newcommand{\TradeoffRows}{%
PostgreSQL & 1 & 7.33 & 156.6 & 5/5 \\
PostgreSQL & 2 & 13.57 & 181.8 & 4/5 \\
PostgreSQL & 5 & 27.90 & 236.7 & 0/5 \\
MongoDB comunitario & 1 & 2.35 & 975.9 & 0/5 \\
MongoDB comunitario & 2 & 4.30 & 1384.1 & 0/5 \\
MongoDB comunitario & 5 & 8.27 & 1719.7 & 0/5 \\
}
\newcommand{\WriteCostRows}{%
PostgreSQL & 5 & 11.391 & 10.458--12.964 \\
MongoDB comunitario & 5 & 0.035 & 0.031--0.037 \\
}

\newcommand{\tablaMotor}[4]{%
\begin{center}
\small
\begin{tabular}{@{}rlrrrrrrr@{}}
\toprule
Pedidos & Consulta & Clientes & Vál./int. & Consultas/s & P50 & P95 [P25, P75] & P99 & $n_{<100}$\\
\midrule
#1
\bottomrule
\end{tabular}
\par\smallskip
\textit{Tabla #3. Rendimiento.} Latencias en ms. P25 y P75 describen la dispersión entre los cinco P95.
\par\medskip
\begin{tabular}{@{}rlrrrrrrr@{}}
\toprule
Pedidos & Consulta & Clientes & CPU (\%) & RSS (MiB) & RAM (MiB) & SoC (°C) & Lect. (MiB) & Escr. (MiB)\\
\midrule
#2
\bottomrule
\end{tabular}
\par\smallskip
\textit{Tabla #4. Recursos.} Medianas por celda de CPU media, RSS máximo, RAM disponible mínima, temperatura máxima y lectura/escritura de la microSD por ventana.
\end{center}
}
\begin{document}
\begin{titlepage}
\thispagestyle{empty}
\centering
\setlength{\parskip}{0pt}
\color{black}
\vspace*{4mm}
{\Large\textbf{EDGE-DB}\par}
\vspace{30mm}
{\fontsize{25}{30}\selectfont\bfseries\color{azul} Agregaciones JSON en una\\Raspberry Pi Zero 2 W\par}
\vspace{7mm}
{\Large PostgreSQL JSONB frente a MongoDB\\comunitario en computación de borde\par}
\vspace{13mm}
{\large Luis Alberto Nafate Hernandez\par}
\vspace{4mm}
Proyecto personal de computación de borde\par
\vfill
{\large\textbf{REPORTE TÉCNICO EXPERIMENTAL}\par}
\vspace{4mm}
{\small 140 corridas principales válidas y 52,998 operaciones cotejadas.\\Resultados, incidentes y limitaciones documentados.\par}
\vspace{7mm}
22 de noviembre de 2026\par
\vspace{2mm}
{\small \textcopyright\ 2026 Luis Alberto Nafate Hernandez. Todos los derechos reservados.\par}
\vspace*{3mm}
\end{titlepage}
\setcounter{page}{2}

\section*{Resumen}
Se evaluaron PostgreSQL 17.12 con JSONB y una compilación comunitaria de MongoDB 7.0.40 en una Raspberry Pi Zero 2 W sin disipador. Se ejecutaron dos consultas equivalentes en cinco tamaños de 1,000 a 50,000 pedidos y, a 10,000 pedidos, con 1, 2 y 5 clientes. Se obtuvieron cinco corridas válidas por combinación de tamaño, consulta y concurrencia para cada motor. Ambos completaron las consultas a 50,000 pedidos con un cliente. PostgreSQL obtuvo mayor throughput en 14 de 14 celdas y menor P95 en 13 de 14. En Q2 a 50,000 pedidos, el P95 mediano fue de 469.6 ms con PostgreSQL y 2163.0 ms con MongoDB. El resultado caracteriza dos configuraciones de extremo a extremo en una placa y una red particulares.

\textbf{Palabras clave.} Computación de borde, JSONB, MongoDB, latencia de cola, reproducibilidad.
\tableofcontents

\clearpage
\section{Plataforma y pregunta de investigación}
La computación de borde permite ejecutar parte del procesamiento cerca del origen de los datos \cite{nist}. Un ensayo previo del autor en Docker Desktop estudió el rendimiento bajo distintos presupuestos de CPU y RAM \cite{antecedente}. EDGE-DB traslada la pregunta a una sola placa física con montaje sin disipador.

\textbf{Pregunta.} ¿Qué compromisos aparecen entre latencia de cola, consultas por segundo y uso de recursos al escalar dos consultas documentales equivalentes en la Raspberry Pi Zero 2 W?

\begin{center}
\small
\begin{tabularx}{.98\linewidth}{@{}lX@{}}
\toprule Aspecto & Configuración observada\\
\midrule
Placa y CPU & Zero 2 W Rev 1.0, cuatro Cortex-A53, frecuencia reportada entre 600 y 1,000 MHz \cite{pi,arm}.\\
Sistema & Raspberry Pi OS Lite de 64 bits, Debian 13.7, kernel \clave{6.18.50+rpt-rpi-v8}.\\
Memoria & 512 MB nominales. Linux reportó 425,116 KiB utilizables y 424,956 KiB de swap con zram.\\
Almacenamiento & microSD ADATA de 64 GB UHS-I, Class 10, V10, con ext4.\\
PostgreSQL & 17.12 ARM64, tabla con clave primaria y columna JSONB, \clave{shared\_buffers} de 128 MiB, \clave{work\_mem} de 4 MiB y JIT activo.\\
MongoDB & 7.0.40 comunitario ARM64, colección con BSON, caché WiredTiger de 0.25 GiB y límite de memoria de 320 MiB.\\
Cliente y red & Laptop Windows por Wi-Fi. PostgreSQL usó \clave{psql} por SSH y MongoDB usó PyMongo por un túnel SSH.\\
Montaje y energía & Sin disipador. Fuente y cable sin identificar. Temperatura ambiente y consumo eléctrico sin medición.\\
\bottomrule
\end{tabularx}
\par\smallskip\textit{Tabla 1.} Plataforma y configuraciones de la ampliación.
\end{center}

\subsection{Factibilidad y versiones}
MongoDB 8.0 oficial para ARM64 requiere ARMv8.2-A o posterior \cite{mongo}. El Cortex-A53 implementa Armv8-A \cite{arm}. Se usó la compilación comunitaria 7.0.40 \cite{community}, con OpenSSL 1.1 dentro de una unidad temporal. La prueba de 1,000 pedidos coincidió con el oráculo antes y después de reiniciar MongoDB. WiredTiger sin atómicas LSE tiene consideraciones de corrección adicionales \cite{wiredtiger}.

\subsection{Pedidos y oráculo}
Los motores usaron los mismos documentos sintéticos. Cada prefijo procede de la semilla 202603 y se identifica por SHA-256. El oráculo Python calcula las filas esperadas de forma independiente.
\begin{center}
\small
\begin{tabular}{@{}rrrrr@{}}
\toprule Pedidos & Artículos & Elegibles en Q2 & JSONL (MiB) & PostgreSQL (MiB)\\
\midrule
\EdgeDatasetRows
\bottomrule
\end{tabular}
\par\smallskip\textit{Tabla 2.} Los cinco conjuntos. El tamaño PostgreSQL incluye la clave primaria.
\end{center}

\clearpage
\section{Método de medición}
\subsection{Consultas y matriz}
\textbf{Q1} expande los ítems de cada pedido, conserva artículos de precio mayor a 50 y calcula suma, promedio y conteo por país y categoría \cite{pgjson,pgfunctions}. \textbf{Q2} filtra pedidos completados entre el 1 de marzo inclusive y el 1 de diciembre exclusivo de 2025. Devuelve hasta 1,000 pares de identificador e importe. Ordena por total descendente y resuelve empates por identificador ascendente. En el piloto de MongoDB con 5,000 pedidos se cotejaron los primeros 1,000 de 3,004 elegibles. A 50,000 pedidos calificaron 30,171.

La matriz abarca siete combinaciones de volumen y concurrencia por motor. Incluye cinco tamaños con un cliente y dos condiciones adicionales con 10,000 pedidos y 2 o 5 clientes. Se ejecutaron Q1 y Q2 en cada combinación, con cinco corridas válidas por celda. Los motores se probaron secuencialmente. Se varió el orden Q1/Q2 con semilla fija. La corrida es la unidad de resumen y todas las repeticiones comparten la placa \cite{benchmark}.

\subsection{Ventanas, telemetría y exclusiones}
Cada sesión principal comenzó con 10 minutos de reposo instrumentado. Los pilotos usaron ventanas de 30 s. La matriz final usó 10 s de calentamiento, 60 s de observación y 120 s de recuperación, extensibles a 300 s. Se prolongó la pausa si la mediana de temperatura de los últimos 15 s superaba en más de 3 °C la referencia de reposo. Se trabajó con caché caliente y se volvieron a cotejar los conjuntos al reanudar.

El cronómetro abarca desde el envío de la consulta hasta que la laptop recibe el resultado completo. Después se verifica la salida. Cada ventana cuenta las operaciones completas, válidas y dentro del plazo. Las operaciones que cruzaron el final quedaron marcadas en el CSV.

La telemetría, cercana a 1 Hz, registra CPU global, RSS agregado del motor, RAM disponible, temperatura del SoC y contadores de microSD. El RSS puede contar páginas compartidas más de una vez. Lecturas y escrituras corresponden al dispositivo completo. El protocolo fijó paradas por 70 °C, menos de 64 MiB de RAM disponible, flags activos, errores de validación o pérdida de telemetría. La reanudación aplicó un máximo de 5 s entre muestras dentro de la ventana. Se conservó el historial de intentos para reconstruir la selección \cite{sandve}.

\subsection{Métricas}
\begin{tabularx}{\linewidth}{@{}lX@{}}
\toprule Métrica & Cálculo e interpretación\\
\midrule
Consultas/s & Operaciones completas, válidas y dentro de la ventana divididas entre 60 s.\\
P50, P95, P99 & Percentiles de latencia de una corrida por rango más cercano. P95 y P99 describen la cola \cite{tail}.\\
Resumen de celda & Medianas de los cinco valores por corrida. Para P95 se incluyen P25 y P75 entre corridas.\\
P99 con pocos datos & Se informa cuántas repeticiones tuvieron menos de 100 operaciones útiles. Su P99 es una observación extrema.\\
\bottomrule
\end{tabularx}

\subsection{Relación con la serie inicial}
La primera serie de PostgreSQL usó ventanas de 5 s y una enmienda a 20 s para Q1 a 50,000 pedidos. Conservó 80 intentos, incluidos seis pilotos y cuatro intentos anteriores a la enmienda. Sus 70 corridas y 2,750 operaciones válidas documentan la evolución del protocolo. La comparación entre los motores corresponde a la ampliación de 60 s.

\clearpage
\section{Resultados consolidados}
\subsection{Análisis exploratorio de la matriz}
De 150 intentos principales, 140 entraron en la comparación. PostgreSQL aportó 70 corridas válidas de 72 intentos y 36,288 operaciones. MongoDB aportó 70 de 78 y 16,710 operaciones. Los intentos excluidos quedan en el anexo. Cada fila de rendimiento indica válidas sobre intentos. La columna $n_{<100}$ cuenta las repeticiones con menos de 100 operaciones completas.

\subsection{Análisis descriptivo de PostgreSQL JSONB}
\tablaMotor{\PgPerformance}{\PgResources}{3}{4}
En Q1 con un cliente, el P95 pasó de 110.9 ms a 1,000 pedidos a 3133.3 ms a 50,000. En Q2 a 10,000 pedidos, el throughput mediano aumentó de 7.33 a 27.90 consultas/s entre uno y cinco clientes y el P95 subió de 156.6 a 236.7 ms. Once corridas quedaron por debajo de 100 operaciones válidas.

\clearpage
\subsection{Análisis descriptivo de MongoDB comunitario}
\tablaMotor{\MongoPerformance}{\MongoResources}{5}{6}
En Q1 a 50,000 pedidos y un cliente, el P95 fue de 4148.1 ms. La misma consulta obtuvo 97.8 ms a 1,000 pedidos. En Q2 a 10,000 pedidos, el throughput aumentó de 2.35 a 8.27 consultas/s entre uno y cinco clientes y el P95 subió de 975.9 a 1719.7 ms. Veinticuatro corridas quedaron por debajo de 100 operaciones válidas.

El CSV por sesión permite rastrear la dispersión y los reintentos detrás de cada fila. Para comparar las cifras de memoria conviene leer juntas la RAM disponible y el RSS, porque PostgreSQL usa varios procesos y MongoDB se ejecutó con un presupuesto de servicio.

\clearpage
\section{Análisis gráfico}
\begin{center}
\begin{tikzpicture}
\begin{axis}[width=.95\linewidth,height=67mm,xmode=log,ymode=log,grid=major,
xtick={1000,5000,10000,25000,50000},xticklabels={1k,5k,10k,25k,50k},
xlabel={Pedidos con un cliente},ylabel={Mediana del P95 (ms)},legend pos=north west,legend columns=2,legend style={font=\small}]
\addplot[azul,thick,mark=*] coordinates {\PgOneVolume}; \addlegendentry{PostgreSQL Q1}
\addplot[verde,thick,mark=square*] coordinates {\MongoOneVolume}; \addlegendentry{MongoDB Q1}
\addplot[azul,thick,dashed,mark=*] coordinates {\PgTwoVolume}; \addlegendentry{PostgreSQL Q2}
\addplot[verde,thick,dashed,mark=square*] coordinates {\MongoTwoVolume}; \addlegendentry{MongoDB Q2}
\end{axis}
\end{tikzpicture}

\textit{Figura 1.} P95 frente al volumen. Ambos ejes usan escala logarítmica. Cada punto resume cinco ventanas de 60 s.
\end{center}
La latencia de Q1 creció con el tamaño en ambos motores. PostgreSQL obtuvo menor P95 en 13 de 14 celdas. En la excepción, Q1 a 1,000 pedidos con un cliente, MongoDB obtuvo 97.8 ms y PostgreSQL 110.9 ms. En Q2 a 50,000 pedidos, la relación entre el P95 observado con MongoDB y el observado con PostgreSQL fue de aproximadamente 4.6.

\begin{center}
\begin{tikzpicture}
\begin{axis}[width=.95\linewidth,height=67mm,xtick={1,2,5},ymin=0,grid=major,
xlabel={Clientes a 10,000 pedidos},ylabel={Mediana de consultas/s},legend pos=north west,legend columns=2,legend style={font=\small}]
\addplot[azul,thick,mark=*] coordinates {\PgOneClients}; \addlegendentry{PostgreSQL Q1}
\addplot[verde,thick,mark=square*] coordinates {\MongoOneClients}; \addlegendentry{MongoDB Q1}
\addplot[azul,thick,dashed,mark=*] coordinates {\PgTwoClients}; \addlegendentry{PostgreSQL Q2}
\addplot[verde,thick,dashed,mark=square*] coordinates {\MongoTwoClients}; \addlegendentry{MongoDB Q2}
\end{axis}
\end{tikzpicture}

\textit{Figura 2.} Throughput frente a concurrencia. El máximo ensayado fue de cinco clientes a 10,000 pedidos.
\end{center}

\clearpage
\subsection{Escritura en microSD y recuperación térmica}
\begin{center}
\begin{tikzpicture}
\begin{axis}[width=.84\linewidth,height=50mm,ybar,bar width=20pt,ymin=0,ymax=260,
symbolic x coords={pg,mongo},xtick=data,xticklabels={PostgreSQL,MongoDB},
ylabel={Escritura por ventana (MiB)},enlarge x limits=.38,ymajorgrids,
nodes near coords,point meta=y,every node near coord/.append style={font=\small}]
\addplot[draw=azul,fill=azul!48] coordinates {(pg,\PgWrites) (mongo,\MongoWrites)};
\end{axis}
\end{tikzpicture}

\textit{Figura 3.} Escritura mediana en la microSD para Q1 con 50,000 pedidos y un cliente. Cinco ventanas válidas de 60 s por configuración.
\end{center}
En esa condición, la microSD registró \PgWrites{} MiB con PostgreSQL y \MongoWrites{} MiB con MongoDB por ventana. PostgreSQL obtuvo además menor P95. Los contadores incluyen toda la actividad de escritura de la tarjeta. El plan de PostgreSQL de la serie inicial mostró un ordenamiento externo, una observación de otra sesión que ayuda a contextualizar el resultado.

\begin{center}
\begin{tikzpicture}
\begin{axis}[width=.94\linewidth,height=53mm,ybar,bar width=14pt,ymin=0,ymax=66,
symbolic x coords={base,extendida,limite},xtick=data,
xticklabels={Recuperó a 120 s,Recuperó con extensión,Alcanzó 300 s sin recuperar},
xticklabel style={font=\footnotesize},ylabel={Pausas evaluables},enlarge x limits=.19,ymajorgrids,
nodes near coords,point meta=y,every node near coord/.append style={font=\footnotesize},legend style={at={(.98,.98)},anchor=north east,font=\footnotesize}]
\addplot[draw=azul,fill=azul!48] coordinates {\PgCooling}; \addlegendentry{PostgreSQL}
\addplot[draw=verde,fill=verde!42] coordinates {\MongoCooling}; \addlegendentry{MongoDB}
\end{axis}
\end{tikzpicture}

\textit{Figura 4.} Pausas evaluables por desenlace. Se considera recuperación cuando la mediana de los últimos 15 s cae dentro del margen de 3 °C sobre el reposo de la sesión.
\end{center}
Se prolongaron \PgCoolingProlonged{} de \PgCoolingEvaluable{} pausas evaluables de PostgreSQL y \MongoCoolingProlonged{} de \MongoCoolingEvaluable{} de MongoDB. Llegaron al límite sin recuperar la banda \PgCoolingUnrecovered{} pausas de PostgreSQL y \MongoCoolingUnrecovered{} de MongoDB. Otras dos pausas de PostgreSQL quedaron incompletas por interrupciones. Las pausas ampliadas forman parte del tiempo experimental observado.

\subsection{Ejemplo de aplicación con un presupuesto de latencia}
\begin{center}
\small
\begin{tabular}{@{}lrrr@{}}
\toprule
Motor & P95 mediano (ms) & Rango de P95 (ms) & Corridas con P95 $\leq 500$ ms\\
\midrule
\BudgetRows
\bottomrule
\end{tabular}
\par\smallskip\textit{Tabla 7.} Q2 con 50,000 pedidos y un cliente. El objetivo de 500 ms es ilustrativo.
\end{center}
La mediana de PostgreSQL quedó dentro del objetivo, aunque una repetición lo superó. Las cinco de MongoDB lo superaron. El recuento muestra cómo cambió el cumplimiento entre ventanas bajo el mismo escenario medido.

\clearpage
\subsection{El rendimiento frente a un límite de latencia}
\begin{center}
\begin{tikzpicture}
\begin{axis}[width=.92\linewidth,height=66mm,ymode=log,xmin=0,xmax=31,ymin=120,ymax=2550,
xlabel={Mediana de consultas/s},ylabel={Mediana del P95 (ms)},grid=major,
ytick={125,200,500,1000,2000},yticklabels={125,200,500,1000,2000},
legend style={at={(.98,.98)},anchor=north east,font=\small}]
\addplot[azul,thick,mark=*,point meta=explicit symbolic,nodes near coords,
every node near coord/.append style={font=\scriptsize,above right}] coordinates {\PgTradeoff};
\addlegendentry{PostgreSQL}
\addplot[verde,thick,mark=square*,point meta=explicit symbolic,nodes near coords,
every node near coord/.append style={font=\scriptsize,above right}] coordinates {\MongoTradeoff};
\addlegendentry{MongoDB}
\addplot[tinta,dashed,thin] coordinates {(0,200) (31,200)};
\addlegendentry{Objetivo de 200 ms}
\end{axis}
\end{tikzpicture}

\textit{Figura 5.} Latencia frente a throughput en Q2 con 10,000 pedidos. Los números junto a los puntos indican clientes. Cada punto combina dos medianas de cinco corridas. El eje vertical es logarítmico.
\end{center}
\begin{center}
\small
\begin{tabular}{@{}lrrrr@{}}
\toprule Motor & Clientes & Consultas/s & P95 (ms) & Corridas con P95 $\leq 200$ ms\\
\midrule
\TradeoffRows
\bottomrule
\end{tabular}
\par\smallskip\textit{Tabla 8.} Medianas por condición y recuento observado de cumplimiento del umbral ilustrativo.
\end{center}
Con dos clientes, PostgreSQL alcanzó 13.57 consultas/s y un P95 mediano de 181.8 ms. Con cinco clientes alcanzó 27.90 consultas/s y 236.7 ms. Dos clientes fue la condición más rápida de las ensayadas que cumplió el objetivo mediano de 200 ms. Lo cumplió en cuatro de cinco corridas, mientras que con un cliente se cumplió en cinco de cinco. Una meta operativa requiere comprobar la variación bajo la carga real de la aplicación.

\subsection{Escritura registrada por consulta completada}
\begin{center}
\small
\begin{tabular}{@{}lrrr@{}}
\toprule Motor & Corridas & Mediana (MiB/consulta) & Rango (MiB/consulta)\\
\midrule
\WriteCostRows
\bottomrule
\end{tabular}
\par\smallskip\textit{Tabla 9.} Q1 a 50,000 pedidos con un cliente. Para cada ventana se dividió la escritura observada en toda la microSD entre las consultas válidas completadas, antes de calcular la mediana y el rango.
\end{center}
El numerador puede incluir actividad del sistema y trabajo de consultas que cruzaron el cierre. Para identificar el origen de las escrituras, una prueba de seguimiento alternaría ventanas de reposo y de Q1 con 50,000 pedidos, un cliente y la misma duración. En cada ventana se registrarían los contadores del dispositivo y, durante la carga, la evidencia de archivos temporales del motor con sus marcas de tiempo.

\clearpage
\section{Análisis diagnóstico y planes de ejecución}
Los planes se obtuvieron fuera de las ventanas cronometradas. PostgreSQL aportó \clave{EXPLAIN ANALYZE} con buffers desde la serie inicial y MongoDB, \clave{executionStats} al cierre de la ampliación. Los tiempos del plan y la latencia de la laptop describen recorridos diferentes \cite{pgexplain}.

\begin{center}
\small
\begin{tabularx}{\linewidth}{@{}lllXr@{}}
\toprule Motor & Consulta & Resultado & Evidencia a 50,000 pedidos & Tiempo (ms)\\
\midrule
PostgreSQL & Q1 & 16 grupos & Escaneo, expansión y agregado. Ordenamiento externo y 3,532 bloques temporales escritos en el plan. & 3179.7\\
PostgreSQL & Q2 & 1,000 filas & Límite y ordenamiento \textit{top-N heapsort}. Cero bloques temporales escritos. & 305.6\\
MongoDB & Q1 & 16 grupos & \clave{COLLSCAN}, expansión a 149,760 ítems y filtro a 106,861. Cero derrames en las etapas que informaron el contador. & 3438\\
MongoDB & Q2 & 1,000 filas & \clave{COLLSCAN}, 50,000 documentos examinados y 30,171 elegibles. Cero derrames en la etapa que informó el contador. & 879\\
\bottomrule
\end{tabularx}
\par\smallskip\textit{Tabla 10.} En MongoDB se indica \clave{executionTimeMillis} de las estadísticas del cursor. En PostgreSQL se indica \clave{Execution Time}.
\end{center}

\subsection{Diagnósticos complementarios}
A 10,000 pedidos, MongoDB completó cinco corridas Q2 con permiso de escritura temporal y cinco sin permiso. Las medianas de P95 fueron 676.8 y 511.1 ms, con 2.63 y 3.17 consultas/s respectivamente. Las condiciones se ejecutaron por bloques. La diferencia puede incluir efectos de sesión y el uso efectivo del disco se verifica en el plan de cada ejecución.

El diagnóstico de PostgreSQL para contrastar 4 y 16 MiB de \clave{work\_mem} se detuvo al agotarse el tiempo límite de consulta en la primera corrida a 4 MiB. El efecto de cambiar a 16 MiB quedó sin estimar. El ordenamiento externo de Q1 procede del plan de la serie inicial.

\section{Discusión y conclusiones}
PostgreSQL procesó más consultas por segundo en las 14 celdas comparables. En las dos configuraciones, aumentar clientes a 10,000 pedidos elevó throughput y latencia de cola. La Q2 de PostgreSQL a cinco clientes ilustra el intercambio. Alcanzó 27.90 consultas/s y 236.7 ms de P95 frente a 7.33 consultas/s y 156.6 ms con un cliente.

Los máximos de temperatura durante ventanas válidas fueron 66.6 °C en PostgreSQL y 68.2 °C en MongoDB. La RAM disponible mínima fue de 158.2 y 150.3 MiB respectivamente. Se observó CPU global cercana a saturación y máscara 0x0 en ambos. Los enfriamientos, la falta de medición ambiental y las sesiones reanudadas limitan la atribución térmica.

\subsection{Amenazas a la validez}
\textbf{Comparabilidad.} Cambian el motor, el cliente, la serialización y el presupuesto de memoria. El cronómetro contiene todos esos componentes. Aislar la contribución del motor requeriría medir por separado el tiempo de servidor y el de transporte, con bloques de carga alternados y RTT instrumentado.

\textbf{Selección e incertidumbre.} Los percentiles describen operaciones aceptadas en ventanas válidas. Los intentos excluidos se informan aparte y la variación entre cinco corridas no mide variación entre placas. El P99 con pocas operaciones es particularmente sensible a un solo valor. Comparar el 2/72 de exclusiones de PostgreSQL con el 8/78 de MongoDB describe lo que ocurrió en el montaje. Los mensajes disponibles no aíslan una tasa de fallo de cada motor.

\textbf{Alcance.} El máximo de 50,000 pedidos corresponde a un cliente y cinco clientes se probaron a 10,000 pedidos. Para ampliar la inferencia se necesitarían otras placas, alternar bloques entre motores, medir la temperatura ambiente y ampliar las ventanas con pocas operaciones.

\clearpage
\appendix
\section{Incidentes y trazabilidad}
\subsection{Exclusiones conservadas}
\begin{center}
\small
\begin{tabularx}{.98\linewidth}{@{}llrX@{}}
\toprule Motor & Incidente & Intentos & Evidencia y tratamiento\\
\midrule
PostgreSQL & Pérdida de telemetría & 2 & Corte en Q1 a 1,000 pedidos y brecha de 6.7 s en Q2 a 10,000 con dos clientes.\\
MongoDB & Brecha de telemetría & 4 & Eventos de 5.2 a 10.1 s en celdas de 10,000, 25,000 y 50,000 pedidos.\\
MongoDB & Tiempo de espera del cliente & 4 & Eventos a 1,000 y 10,000 pedidos. El CSV conserva la condición y el mensaje.\\
\bottomrule
\end{tabularx}
\par\smallskip\textit{Tabla A1.} Diez exclusiones de la matriz principal. Las incidencias fuera de las corridas se conservan por sesión.
\end{center}
Una interrupción de SSH dejó activo un servicio temporal de MongoDB. La recuperación comprobó el fin del servicio y el retorno de PostgreSQL antes de reanudar. Los incidentes de conexión y telemetría forman parte de la evidencia sobre el montaje.

\subsection{Cronología de una exclusión y su reintento}
\begin{center}
\small
\begin{tabularx}{.98\linewidth}{@{}lX@{}}
\toprule 19 de noviembre, hora local & Q2 con 10,000 pedidos, dos clientes, repetición 5 de PostgreSQL\\
\midrule
22:18:53 & Inició el intento V02. La última muestra dentro de la ventana se recibió a los 53.26 s. Quedaron 6.74 s hasta el cierre, por encima del umbral de 5 s.\\
22:20:13 & El supervisor cerró la sesión y conservó las 560 operaciones completas del intento. Esa ventana quedó excluida. Las 38 corridas válidas previas permanecieron disponibles para reanudar.\\
22:32:13 & El intento V03 repitió la misma condición con el mismo identificador de corrida. Completó 835 operaciones válidas en una ventana aceptada y aportó la repetición a la matriz final.\\
\bottomrule
\end{tabularx}
\par\smallskip\textit{Tabla A2.} El identificador de sesión distingue el primer intento de su reintento. La marca corresponde al inicio de la corrida o al cierre documentado de la sesión.
\end{center}
La regla de exclusión comprueba la cobertura de telemetría. El reintento restituye la réplica planificada y el registro conserva el intento anterior para reconstruir la selección.

\subsection{Verificación eléctrica}
En un ensayo separado se guardaron 64 lecturas del voltaje interno VC4. Las medianas fueron 1.2500 V en reposo, 1.2625 V en Q1 a 10,000 pedidos y 1.2500 V durante recuperación. El firmware rechazó la consulta del canal de entrada de 5 V. La familia Zero tiene limitaciones de aviso de bajo voltaje \cite{firmware,power}. El voltaje de entrada y la energía quedaron sin medir.

\subsection{Disponibilidad de datos y código}
\begin{tabularx}{\linewidth}{@{}lX@{}}
\toprule Ruta en el repositorio & Evidencia\\
\midrule
\clave{dataset/} & Generador, semilla y manifiesto SHA-256.\\
\clave{workloads/} & SQL y pipelines de las dos consultas.\\
\clave{results/raw/} & Estados de sesión, CSV por corrida y operación, telemetría, enfriamiento, fallos y planes.\\
\clave{results/processed/} & Resúmenes por sesión y celda.\\
\clave{reporte/} & Fuente, tablas generadas y comando de reconstrucción que coteja las mediciones por operación.\\
\clave{protocol/} & Decisiones, enmiendas, bitácora, prompts y auditoría de IA.\\
\bottomrule
\end{tabularx}

\textbf{Revisión.} El repositorio del proyecto conserva los datos y la ruta de reconstrucción. Codex asistió en código, ejecución, análisis y redacción. La variante comunitaria y los límites del montaje se documentan como parte del alcance experimental.

\clearpage
\phantomsection
\addcontentsline{toc}{section}{Referencias}
\begin{thebibliography}{99}
\small\raggedright
\bibitem{nist} M. Iorga y otros. \textit{Fog Computing Conceptual Model}. NIST SP 500-325, 2018. \href{https://doi.org/10.6028/NIST.SP.500-325}{DOI: 10.6028/NIST.SP.500-325}.
\bibitem{pgjson} PostgreSQL Global Development Group. \textit{PostgreSQL 17: JSON Types}. \href{https://www.postgresql.org/docs/17/datatype-json.html}{Documentación oficial}.
\bibitem{pi} Raspberry Pi Ltd. \textit{Zero models}. \href{https://www.raspberrypi.com/documentation/computers/zero.html}{Especificaciones oficiales}.
\bibitem{mongo} MongoDB Inc. \textit{MongoDB 8.0: Production Notes}, apartado ARM64. \href{https://www.mongodb.com/docs/v8.0/administration/production-notes/#arm64}{Requisitos oficiales}.
\bibitem{arm} Arm. \textit{Cortex-A53 MPCore Processor Technical Reference Manual}, sección 1.2.1. \href{https://documentation-service.arm.com/static/6040c321ee937942ba301626}{Manual de arquitectura}.
\bibitem{pgfunctions} PostgreSQL Global Development Group. \textit{PostgreSQL 17: JSON Functions and Operators}. \href{https://www.postgresql.org/docs/17/functions-json.html}{Funciones de JSONB}.
\bibitem{tail} J. Dean y L. A. Barroso. \textit{The Tail at Scale}. Communications of the ACM 56(2), 2013. \href{https://research.google/pubs/the-tail-at-scale/}{Publicación de los autores}.
\bibitem{pgexplain} PostgreSQL Global Development Group. \textit{PostgreSQL 17: EXPLAIN}. \href{https://www.postgresql.org/docs/17/sql-explain.html}{Referencia de SQL}.
\bibitem{benchmark} T. Kalibera y R. E. Jones. \textit{Rigorous Benchmarking in Reasonable Time}. ISMM, 2013. \href{https://doi.org/10.1145/2464157.2464160}{DOI: 10.1145/2464157.2464160}.
\bibitem{firmware} Raspberry Pi Ltd. \textit{Raspberry Pi OS: vcgencmd}. \href{https://www.raspberrypi.com/documentation/computers/os.html#measure_volts-block}{Voltajes internos y comandos de firmware}.
\bibitem{power} Raspberry Pi Ltd. \textit{Raspberry Pi computer hardware: Power supply warnings}. \href{https://www.raspberrypi.com/documentation/computers/raspberry-pi.html#power-supply-warnings}{Detección de bajo voltaje por familia}.
\bibitem{community} themattman. \textit{r7.0.40 Unofficial Raspberry Pi Binaries}. \href{https://github.com/themattman/mongodb-raspberrypi-binaries/releases/tag/r7.0.40-rpi-unofficial}{Distribución comunitaria y SHA-256}.
\bibitem{wiredtiger} MongoDB. \textit{Investigate changes to WiredTiger for non-LSE ARM platforms}. \href{https://jira.mongodb.org/browse/WT-9256}{Registro técnico WT-9256}.
\bibitem{antecedente} L. A. Nafate Hernandez. \textit{Medición de métricas en bases documentales: MongoDB frente a PostgreSQL JSONB bajo contención de CPU y memoria}. Informe técnico previo del autor, 2026.
\bibitem{sandve} G. K. Sandve y otros. \textit{Ten Simple Rules for Reproducible Computational Research}. PLOS Computational Biology, 2013. \href{https://doi.org/10.1371/journal.pcbi.1003285}{DOI}.
\bibitem{mensh} B. Mensh y K. Kording. \textit{Ten Simple Rules for Structuring Papers}. PLOS Computational Biology, 2017. \href{https://doi.org/10.1371/journal.pcbi.1005619}{DOI}.
\end{thebibliography}
\end{document}
